% !TEX root = ../main.tex
\subsection*{Definition 5.1 --- $L^1$-Space}%
\manualtoc{subsection}{\texorpdfstring{Definition 5.1 --- $L^1$-Space}{Definition 5.1 --- L1-Space}}{studiabile-nav-2006}%

Let $(S,\Sigma,\mu)$ be a measure space. We denote by $L^1(S,\Sigma,\mu)$ the collection of measurable functions $f\in m\Sigma$ such that
\[
\mu(|f|)=\int_S f^+\,d\mu+\int_S f^-\,d\mu<\infty.
\]
Recall from HW~49 that
\[
f=f^+-f^-,\qquad |f|=f^++f^-,
\]
where $f^+,f^-\in(m\Sigma)^+$. The equality in the definition follows from Lemma~5.6.

In particular, if $f\in L^1(S,\Sigma,\mu)$, both $\mu(f^+)$ and $\mu(f^-)$ are finite. Therefore,
\[
\mu(f)=\mu(f^+)-\mu(f^-)\in\mathbb R.
\]
Such functions are called Lebesgue integrable.

We work with measurable representatives. When $L^1(S,\Sigma,\mu)$ is viewed as a normed space, functions equal $\mu$-almost everywhere are identified; the norm of their common equivalence class is $\|f\|_1:=\mu(|f|)$. Before this identification, the same formula defines a seminorm on the collection of integrable functions.

\subsection*{HW 50 --- Absolute Value of the Integral}%
\manualtoc{subsection}{\texorpdfstring{HW 50 --- Absolute Value of the Integral}{HW 50 --- Absolute Value of the Integral}}{studiabile-nav-2007}%

\studyheading{Statement.} For $f\in L^1(S,\Sigma,\mu)$, prove that
\[
|\mu(f)|\le\mu(|f|).
\]
\par\noindent\studyheading{Expanded Proof.}\par\nobreak

Since $f\in L^1(S,\Sigma,\mu)$, we have
\[
\mu(f^+)+\mu(f^-)=\mu(|f|)<\infty.
\]
Thus both integrals are finite. By the definition of the signed Lebesgue integral,
\[
\begin{aligned}
|\mu(f)|&=|\mu(f^+)-\mu(f^-)| \le\mu(f^+)+\mu(f^-) =\mu(|f|).
\end{aligned}
\]
The inequality follows from the triangle inequality for real numbers, since $\mu(f^+),\mu(f^-)\ge0$, and the last equality follows from Definition~5.1.

Therefore, the statement is proved. 

\subsection*{HW 51 --- Closure of $L^1$ under Linear Combinations}%
\manualtoc{subsection}{\texorpdfstring{HW 51 --- Closure of $L^1$ under Linear Combinations}{HW 51 --- Closure of L1 under Linear Combinations}}{studiabile-nav-2008}%

\studyheading{Statement.} Let $f,g\in L^1(S,\Sigma,\mu)$ and $\alpha,\beta\in\mathbb R$. Prove that
\[
\alpha f+\beta g\in L^1(S,\Sigma,\mu).
\]
\par\noindent\studyheading{Expanded Proof.}\par\nobreak

Set $u:=\alpha f+\beta g$. By Lemma~3.4, $u\in m\Sigma$, since $f,g$ are measurable and $\alpha,\beta\in\mathbb R$.

For every $s\in S$, the triangle inequality gives
\[
|u(s)|=|\alpha f(s)+\beta g(s)|
\le|\alpha|\,|f(s)|+|\beta|\,|g(s)|.
\]
By monotonicity (HW~45) and nonnegative linearity (Lemma~5.6),
\[
\begin{aligned}
\mu(|u|)&\le\mu(|\alpha|\,|f|+|\beta|\,|g|) =|\alpha|\,\mu(|f|)+|\beta|\,\mu(|g|) <\infty.
\end{aligned}
\]
The last inequality follows from $f,g\in L^1(S,\Sigma,\mu)$ and the finiteness of $\alpha,\beta$. Therefore $u\in L^1(S,\Sigma,\mu)$.

Therefore, the statement is proved. 

\subsection*{Linearity of the Integral on $L^1$}%
\manualtoc{subsection}{\texorpdfstring{Linearity of the Integral on $L^1$}{Linearity of the Integral on L1}}{studiabile-nav-2009}%

Linearity extends to signed integrable functions. If $f,g\in L^1(S,\Sigma,\mu)$ and $\alpha,\beta\in\mathbb R$, then
\[
\int_S(\alpha f+\beta g)\,d\mu
=\alpha\int_S f\,d\mu+\beta\int_S g\,d\mu.
\]
\par\noindent\studyheading{Expanded Proof.}\par\nobreak

By HW~51, $u:=\alpha f+\beta g\in L^1(S,\Sigma,\mu)$. Define the positive and negative parts of the scalars by
\[
\alpha^+=\max\{\alpha,0\},\quad\alpha^-=\max\{-\alpha,0\},
\qquad
\beta^+=\max\{\beta,0\},\quad\beta^-=\max\{-\beta,0\}.
\]
Consider the nonnegative measurable functions
\[
\begin{aligned}
P&:=\alpha^+f^++\alpha^-f^-+\beta^+g^++\beta^-g^-,\\
Q&:=\alpha^+f^-+\alpha^-f^++\beta^+g^-+\beta^-g^+.
\end{aligned}
\]
Since $f,g\in L^1$ and all coefficients are nonnegative, Lemma~5.6 gives $\mu(P),\mu(Q)<\infty$. Moreover,
\[
P-Q=(\alpha^+-\alpha^-)(f^+-f^-)
+(\beta^+-\beta^-)(g^+-g^-)
=\alpha f+\beta g=u.
\]
Together with $u=u^+-u^-$, this yields the pointwise identity
\[
u^++Q=u^-+P.
\]
By Lemma~5.6,
\[
\mu(u^+)+\mu(Q)=\mu(u^-)+\mu(P).
\]
All terms are finite, since $u\in L^1$ and $\mu(P),\mu(Q)<\infty$. Thus,
\[
\begin{aligned}
\mu(\alpha f+\beta g)&=\mu(u^+)-\mu(u^-) =\mu(P)-\mu(Q)\\
&=(\alpha^+-\alpha^-)(\mu(f^+)-\mu(f^-)) \quad +(\beta^+-\beta^-)(\mu(g^+)-\mu(g^-)) =\alpha\mu(f)+\beta\mu(g).
\end{aligned}
\]
The second equality follows from the preceding integrated identity; the third uses Lemma~5.6 and rearrangement of finite quantities, while the last follows from the signed-integral definition.

Therefore, the statement is proved. 

We will often need to interchange limits and integrals. One result of this kind is MON; the next theorem provides another important tool for doing so.
